\defn{Pairwise sequence alignment} is a computational step commonly required in bioinformatics pipelines~\citep{alser2022frommolecules}, such as in \defn{read mapping}~\cite{alser2020accelerating} and \defn{de-novo assembly}~\citep{li2011denovo}. We formulate the problem as (1)~finding the \defn{edit distance} between two sequences~\citep{levenshtein1966binary} and (2)~determining the sequence of corresponding edits. Efficient algorithms for solving this problem optimally are based on \defn{dynamic programming (DP)}, such as the Smith-Waterman-Gotoh algorithm~\citep{smith1981identification,gotoh1982improved}, and have a runtime that grows quadratically with sequence length~\citep{alser2020technology}. \citep{asm_lower_bound}~proves no strongly subquadratic time solutions can exist, provided the strong exponential time hypothesis~\citep{impagliazzo2001complexity} holds. Hence, recent works focus on approaches such as pre-alignment filtering~(e.g.,~\citep{xin2013fasthash,xin2015shiftedhammingdistance,alser2019shouji,alser2020sneakysnake,singh2021fpga,mansouri2022genstore}), constant factor algorithmic speedups~(e.g.,~\citep{edlib, ksw2_a, ksw2_b, wfa_algorithm}), GPU-based acceleration (e.g., \citep{gasal2,ahmed2020darwin_gpu,de2016cudalign,liu2013cudasw++,awan2020adept}), FPGA-based acceleration~(e.g.,~\citep{fei2018fpgasw, hoffmann2016using, benkrid2009highly}) or using specialized hardware accelerators~(e.g., \citep{ turakhia2019darwin,turakhia2019darwinwga,fujiki2018genax,fujiki2020seedex,cali2020genasm,cali2022segram}).

We observe that GenASM~\citep{cali2020genasm}, a recent state-of-the-art sequence alignment algorithm, has a large space for improvement. %
GenASM uses only cheap bitwise operations and breaks the lower complexity bound of pairwise sequence alignment through its powerful \defn{windowing heuristic}. \citep{cali2020genasm} has already proven the effectiveness of the GenASM algorithm and its accelerator implementation, thus we are motivated to further improve the GenASM algorithm and explore its potential on commodity hardware.

We identify \textbf{three inefficiencies} in the GenASM algorithm: (1)~it has a \emph{large memory footprint} due to the large size of the dynamic programming (DP) table, (2)~it has a \emph{high amount of data movement} between registers and memory due to frequent accesses to the DP table, and (3)~it does some \emph{unnecessary} work by calculating DP cells that are not useful for finding the final result.
The three inefficiencies negatively impact both (1)~software implementations running on commodity hardware (e.g., CPUs or GPUs) and (2)~custom hardware (e.g., ASIC) implementations.

\emph{Software implementations} on commodity hardware typically cannot
fit all the data into fast on-chip memories (e.g., L1, scratchpad memory) due to the large memory footprint. This increases the latency and limits the bandwidth with which the DP table can be accessed. The high amount of data movement puts high pressure on this bandwidth, limiting performance.

In contrast, \emph{custom hardware implementations} can use arbitrarily large amounts of on-chip memory, but such a large on-chip memory with the high bandwidth requirement is costly. For example, the hardware accelerator described in~\citep{cali2020genasm} requires 76\% and 54\% of the total chip area and power consumption
for the on-chip memory that stores the DP table.

The unnecessary work stems from computing cells that do not contain useful information for finding the final result. This applies to at least 25\% of cells on average for uncorrelated string pairs, and more for correlated string pairs, as we show in \secref{early_termination}. Doing unnecessary work affects software and hardware implementations equally because both could use the wasted time to do useful work instead.


Our \textbf{goal}
is to develop a fast and memory-frugal alignment algorithm by addressing the inefficiencies in the GenASM algorithm, and demonstrate its benefits with high-performance CPU and GPU implementations.


To this end we propose \defn\toolname\footnote{\emph{\toolname} is a name for miserly or frugal fictional characters (e.g.,~\citep{dickens1843achristmascarol}), similar to how our proposed algorithm aims to be as resource-efficient as possible.}, which includes improvements to the GenASM algorithm based on three \textbf{key~ideas}:
\begin{enumerate}
    \item The DP table can be \emph{compressed} by storing only the bitwise AND of multiple values (\secref{SENE}). The required regions of the DP table can then be decompressed on-demand during traceback with a small computational overhead.
    \item Part of the DP table \emph{does not need to be stored} because the traceback operation cannot reach these entries (\secref{DENT}).
    \item Part of the DP table can opportunistically be \emph{excluded from calculation} if previous rows of the DP table already contain the information needed for finding the final result (\secref{early_termination}).
\end{enumerate}
These improvements (1)~reduce the number of accesses to GenASM's DP table, (2)~reduce the memory footprint of the DP table, and (3)~eliminate unnecessary work. %

We experimentally demonstrate that our improvements yield significant benefits across multiple computing platforms and multiple baseline sequence alignment methods. The CPU version of \toolname achieves a \kswcpuspeedup, \edlibcpuspeedup, and \cpuoptspeedup speedup over CPU-based implementations of KSW2~\citep{ksw2_a,ksw2_b}, Edlib~\citep{edlib}, GenASM, respectively. The GPU version of \toolname achieves a \cpugpuspeedup, \kswgpuspeedup, \edlibgpuspeedup, \darwingpuspeedup, and \gpuoptspeedup speedup over CPU-based implementations of \toolname, KSW2, and Edlib, and GPU-based implementations of Darwin-GPU~\citep{ahmed2020darwin_gpu} and GenASM, respectively. We analytically estimate an ASIC implementation of \toolname to use 3.6\x less chip area and consume 2.1\x less power compared to the prior state-of-the-art ASIC implementation of GenASM~\citep{cali2020genasm} while maintaining the same throughput.

The \textbf{contributions} of this paper are as follows:
\begin{itemize}
    \item We develop three novel algorithmic improvements that are applicable to software and custom hardware implementations of \toolname, collectively reducing the memory footprint by 24\x, the number of memory accesses by 12\x, and the number of entries of the DP table calculated by at least 25\% on average compared to GenASM.
    \item We experimentally demonstrate the significant throughput (i.e., alignments per second) increase of our improvements for CPU and GPU implementations of \toolname.
    \item We analytically estimate that an ASIC implementation of \toolname\ significantly reduces the chip area and power consumption compared to the prior state-of-the-art ASIC implementation of GenASM.
    \item We open-source all code, including high-performance CPU and GPU implementations of \toolname, which can be readily used as a sequence alignment library, and all evaluation scripts.
    \item We systematically analyze the throughput and accuracy behavior of GenASM and \toolname across a range of configurations based on real and simulated datasets for long and short reads. As the best configuration of \toolname depends on the computing platform, we make several observations that can help guide future implementations of \toolname.
\end{itemize}
